\section{ZMP 挙動から見た安定性の違い}
\label{sec:discussion_zmp}

ZMP の挙動に着目すると，各 Case に共通して，DS 期間においては ZMP が支持脚側の外側境界付近に張り付く傾向が見られた．このことから，DS 期間における ZMP の位置そのものは，提案手法によって大きく変化していないことが分かる．

しかしながら．DS から SS へ切り替わった直後の挙動に関しては，Case4における提案手法適用時には，SS 期間においても ZMP が支持脚足裏領域内に留まり，発散することなく安定した挙動を示した．

このことは，提案手法が DS--SS 遷移後の ZMP の発散を抑制する効果を持つことを示しており，歩行安定性の向上において重要な役割を果たしていると考えられる．
